\section{Introduction}

The 2030 congressional reapportionment will be the most geographically significant
since the 1990 cycle, when the Sun Belt's post-WWII growth first produced major
reapportionment from the Northeast and Midwest. The 2020-2030 shift is projected to
be even larger in seat magnitude (13+ seats changing states) because the Sun Belt
growth rate has accelerated relative to the pre-2000 period.

This paper examines the algorithmic implications of the projected changes,
focusing on two categories of states: gainers (TX, FL, AZ, GA, CO, NC, WA, OR)
whose seat allocations will require new bisection tree designs, and losers (NY, CA,
IL, OH, MI, PA, WV, MN, RI, MT) whose seat reductions will require district mergers
or major boundary redesigns.

The central algorithmic finding: GeoSection (\texttt{--structure ratio-optimal})
is the recommended structure for the two largest gainers (TX k=41, FL k=30) because
it handles non-standard prime factorizations naturally. ApportionRegions (prime-factor
tree) works correctly for both via its binary fallback, but GeoSection produces
more geographically coherent results in fast-growth states where the natural
urban-suburban-rural gradient aligns with ratio-optimal bisection.
